\begin{center}
\hypertarget{native-057-article039}{}\pdfbookmark[2]{39.}{article-039}39.
\end{center}
Dimidiam fere partem laboris quem illa tentamina requirerent abscindi posse statim obvium est, si tabula ita adornata habeatur, ex qua $\log B$ immediate argumento $A$ desumere liceat. Tres tunc superessent operationes; prima indirecta, puta determinatio ipsius $A$, ut aequationi 1) art. 37 satisfiat; secunda, determinatio ipsius $E$ ex $A$ et $B$, quae fit directe vel per aequationem $E=2B(A^{\frac{1}{2}}+\tfrac{1}{15}A^{\frac{3}{2}})$, vel per hanc $\sin E=2B(A^{\frac{1}{2}}-\tfrac{3}{5}A^{\frac{3}{2}})$; tertia, determinatio ipsius $v$ ex $E$ per aequ. VII art. 8. Operationem primam ad algorithmum expeditum et a tentaminibus vagis liberum reducemus; secundam et tertiam vero in unicam contrahemus, tabulae nostrae quantitatem novam $C$ inserendo, quo pacto ipsa $E$ omnino opus non habebimus, simulque pro radio vectore formulam elegantem et commodam nanciscemur. Quae singula ordine suo iam persequemur.

Primo aequationem 1) ita transformabimus, ut tabulam \textsc{Barker}ianam ad eius solutionem adhibere liceat. Statuemus ad hunc finem $A^{\frac{1}{2}}=\operatorname{tang}\tfrac{1}{2}w\,\mathord{.}\,\BareRoot{\frac{5-5e}{1+9e}}$, unde fit
\[
75\operatorname{tang}\tfrac{1}{2}w+25\operatorname{tang}\tfrac{1}{2}w^3=\frac{75kt\BareRoot{(\tfrac{1}{5}+\tfrac{9}{5}e)}}{2Bq^{\frac{3}{2}}}=\frac{\alpha t}{B},
\]
designando constantem $\frac{75k\BareRoot{(\tfrac{1}{5}+\tfrac{9}{5}e)}}{2q^{\frac{3}{2}}}$ per $\alpha$\HistoricalNoteEight{Handschriftliche Bemerkung:] Wenn man \textsc{Burckhardts} Tafel [Connaissance des Temps 1818; s. auch Bd. VI. S. 596] braucht, $\frac{\BareRoot{(\tfrac{1}{10}+\tfrac{9}{10}e)}}{q^{\frac{3}{2}}}=\alpha$.}. Si itaque $B$ esset cognita, $w$ illico \SourceLeaf{058}e tabula \textsc{Barker}iana desumi posset, ubi est anomalia vera, cui respondet motus medius $\frac{\alpha t}{B}$; ex $w$ derivabitur $A$ per formulam $A=\beta\operatorname{tang}\tfrac{1}{2}w^2$, designando constantem $\frac{5-5e}{1+9e}$ per $\beta$. Iam etsi $B$ demum ex $A$ per tabulam nostram auxiliarem innotescat, tamen propter perparvam ipsius ab unitate differentiam praevidere licet, $w$ et $A$ levi tantum errore affectas provenire posse, si ab initio divisor $B$ omnino negligatur. Determinabimus itaque primo, levi tantum calamo, $w$ et $A$, statuendo $B=1$; cum valore approximato ipsius $A$ e tabula nostra auxiliari inveniemus ipsam $B$, cum qua eundem calculum exactius repetemus; plerumque respondebit valori sic correcto ipsius $A$ prorsus idem valor ipsius $B$, qui ex approximato inventus erat, ita ut nova operationis repetitio superflua sit, talibus casibus exceptis, ubi valor ipsius $E$ iam valde considerabilis fuerit. Ceterum vix opus erit monere, si forte iam ab initio valor ipsius $B$ quomodocunque approximatus aliunde innotuerit (quod semper fiet, quoties e pluribus locis, haud multum ab invicem distantibus, computandis unus aut alter iam sunt absoluti), praestare, hoc statim in prima approximatione uti: hoc modo calculator scitus saepissime ne una quidem calculi repetitione opus habebit. Hanc celerrimam approximationem inde assecuti sumus, quod $B$ ab $1$ differentia ordinis quarti tantum distat, in coëfficientem perparvum numericum insuper multiplicata, quod commodum praeparatum esse iam perspicietur per introductionem quantitatum $E-\sin E$, $\tfrac{9}{10}E+\tfrac{1}{10}\sin E$ loco ipsarum $E$, $\sin E$.
